\section{Stability for concave equations}
\label{sec:concave-strict-gap-stability}\index{Stability!concave equations}

For a concave equation, an admissible reference with a strict asymptotic gap replaces the determinant bounds used above. It bounds positive matrix perturbations at the reference. An auxiliary Monge--Amp\`ere equation then turns a small superlevel set of the difference of two solutions into a contradiction. The resulting stability estimate remains uniform when their positive right-hand sides approach zero.

Return to a connected compact complex manifold $X$ of dimension $n\geq2$ with a fixed Hermitian form $\omega$. Fix a smooth real $(1,1)$-form $\chi$. For a smooth real function $w$, put
\[
 A_w=\omega^{-1}(\chi+\ddc w),\qquad
 F(A)=f(\lambda(A)).
\]
Let $\Gamma$ be an open convex symmetric cone with $\Gamma_n\subseteq\Gamma\subseteq\Gamma_1$, where $\Gamma_n$ is the positive orthant and $\Gamma_1=\{\sum_i\lambda_i>0\}$. Assume that $f\in C^\infty(\Gamma)\cap C^0(\overline\Gamma)$ is symmetric, concave, positive on $\Gamma$, zero on $\partial\Gamma$, and satisfies
\[
 f_i>0,\qquad \lim_{t\to\infty}f(t\lambda)=+\infty
 \quad(\lambda\in\Gamma).
\]
Call $w$ admissible when $\lambda(A_w)\in\Gamma$ everywhere. For $\mu\in\Gamma$, define
\[
 f_{\infty,i}(\mu)=\lim_{t\to\infty}f(\mu+te_i),\qquad
 f_\infty(\mu)=\min_i f_{\infty,i}(\mu),
\]
allowing $+\infty$. Fix one smooth admissible $\underline u$ and positive constants $H_0,\delta_0$. Consider smooth admissible solutions of
\[
 F(A_u)=h,\qquad 0<h\leq H_0,\qquad
 f_\infty(\lambda(A_{\underline u}(x)))\geq h(x)+\delta_0.
\]
Use $\inf_Xu=0$ to fix additive constants. All constants in this section may depend on $X,\omega,\chi,f,\Gamma,\underline u,H_0,\delta_0$. Norms use $\omega^n$, and $\langle a\rangle_\omega=(\int_X\omega^n)^{-1}\int_Xa\omega^n$ denotes the average.

\Needspace{10\baselineskip}
\begin{lemma}\label{lem:strict-gap-matrix-bound}
There are a positive continuous $H$ and constants $r,R,\delta>0$ such that $h\leq H$ for every right-hand side in the family,
\[
 \mu-2r\mathbf1\in\Gamma,\qquad
 f_\infty(\mu-2r\mathbf1)\geq H+3\delta,
 \qquad \mu=\lambda(A_{\underline u}),
\]
and every positive semidefinite Hermitian endomorphism $P$ satisfies
\[
 F(A_{\underline u}-r\operatorname{Id}+P)\leq H+2\delta
 \quad\Longrightarrow\quad P\leq R\operatorname{Id}.
\]
\end{lemma}
\begin{proof}
Each $f_{\infty,i}$ is an increasing limit of positive concave functions. If it is infinite at one point and finite at another interior point, extend the segment slightly beyond the latter point. The latter point is an interior point of this longer segment, so concavity forces its value to be infinite, a contradiction. Thus it is finite everywhere and continuous, or infinite everywhere. Symmetry gives the same alternative for every $i$.

Take $\delta=\frac18\min\{\delta_0,\min_Xf(\mu)\}>0$. In the finite case set $H=\min\{H_0,f_\infty(\mu)-4\delta\}$; otherwise set $H=H_0$. Since $f_\infty(\mu)\geq f(\mu)\geq8\delta$, this function is positive. The family hypothesis gives $h\leq H$ and $f_\infty(\mu)\geq H+4\delta$. Compactness of $\mu(X)\Subset\Gamma$ and continuity, or divergence, of the asymptotic slopes give a sufficiently small common $r$ with the displayed gap.

If the matrix bound failed, there would be $x_j\to x_*$ and $P_j\geq0$ with unbounded norm and $F(A_{\underline u}(x_j)-r\operatorname{Id}+P_j)\leq H(x_j)+2\delta$. Their ordered eigenvalues $\lambda^{(j)}$ satisfy $\lambda_i^{(j)}\geq\mu_i(x_j)-r$ by eigenvalue monotonicity. Their sum tends to infinity, so some component $k$ tends to infinity along a subsequence. For every fixed $t>0$, eventually
\[
 \lambda^{(j)}\geq\mu(x_*)-2r\mathbf1+te_k.
\]
Ellipticity, followed by $t\to\infty$, gives a lower limit at least $H(x_*)+3\delta$, contradicting the upper bound. The cone is preserved by positive semidefinite matrix addition, so all these comparisons are admissible.
\end{proof}

\Needspace{9\baselineskip}
\begin{proposition}\label{prop:strict-gap-potential-bound}
Every solution satisfies $0\leq u\leq C$.
\end{proposition}
\begin{proof}
Put $\theta=\chi+\ddc\underline u$ and $w=u-\underline u-\sup_X(u-\underline u)\leq0$. Since $\Gamma\subseteq\Gamma_1$, $\operatorname{tr}_\omega\ddc w\geq-\operatorname{tr}_\omega\theta\geq-C$. The fixed Green-kernel construction in the proof of \cref{prop:stw-czero-reconstructed}, applied to $\operatorname{tr}_\omega\ddc$, bounds $\|w\|_{L^1}$ uniformly.

Center a uniformly sized holomorphic coordinate ball at a minimum point of $w$ and rescale it to the real unit ball. For a fixed sufficiently small $\varepsilon>0$, set $q=w+\varepsilon|x|^2$ and $m=\min_Xw$. The real contact-set construction of \eqref{eq:stw-contact-area} gives a compact set $E$ on which
\[
 D^2q\geq0,\qquad q\leq m+\varepsilon/2,
 \qquad |B_{\varepsilon/2}|\leq\int_E\det D^2q\,\mathrm dx.
\]
On $E$, $\ddc w\geq-C\varepsilon\omega$. Choose $\varepsilon$ so that $C\varepsilon<r$. Then $P=\omega^{-1}(\ddc w+r\omega)\geq0$ and
\[
 F(A_{\underline u}-r\operatorname{Id}+P)=F(A_u)=h\leq H.
\]
\Cref{lem:strict-gap-matrix-bound} bounds $P$ and hence the complex trace of $q$. The nonnegative real Hessian has a bounded trace and therefore a bounded determinant. The contact inequality implies $|E|\geq c>0$. Since $w\leq m+\varepsilon/2$ there, the $L^1$ bound gives $c(-m-\varepsilon/2)\leq C$. This bounds the oscillation of $w$, hence that of $u$, proving the assertion.
\end{proof}

\Needspace{11\baselineskip}
\begin{theorem}\label{thm:concave-linear-stability}\index{Stability!linear density and L1 estimate@linear density and $L^1$ estimate}
Two solutions $u_1,u_2$ with the same fixed data and right-hand sides $h_1,h_2$ satisfy
\[
 \|u_1-u_2\|_\infty\leq C\bigl(\|h_1-h_2\|_\infty+\|u_1-u_2\|_{L^1}\bigr).
\]
The constant requires neither a positive common lower bound for $h_i$ nor bounds for their derivatives.
\end{theorem}
\begin{proof}
Set $a=u_1-\underline u$, $b=u_2-\underline u$ and $\theta=\chi+\ddc\underline u$. By \cref{prop:strict-gap-potential-bound}, $|a|,|b|\leq C_0$ for a fixed $C_0\geq1$. It suffices to bound $s=\max_X(b-a)>0$; the opposite difference follows by interchange.

Fix $r,R,\delta$ from \cref{lem:strict-gap-matrix-bound} and put $\alpha=r\omega$. For every smooth positive $\rho$ with
\[
 \rho\geq\tfrac12,\qquad
 \langle\rho\rangle_\alpha=1,\qquad
 \langle\rho^2\rangle_\alpha\leq4,
\]
\cref{thm:tw-smooth} gives a smooth solution
\[
 (\alpha+\ddc\psi)^n=c\rho\alpha^n,\qquad
 \alpha+\ddc\psi>0,\qquad\sup_X\psi=0.
\]
At a maximum of $\psi$, $c\rho\leq1$, so $c\leq2$. Apply \cref{thm:lp-potential-estimate} with $p=2$ and background $\alpha$. Its potential and lower-density bounds give fixed $c_0,C_\psi>0$ with
\[
 c\geq c_0,\qquad -C_\psi\leq\psi\leq0.
\]
These constants depend only on the three displayed density bounds, before any further cutoff height is chosen.

Choose $\kappa>0$ so that $2C_0\kappa\leq1/4$, $\kappa C_0\leq1/16$ and $\kappa C_\psi/2\leq1/16$. Next choose $M\geq1$ with $2^{-n}c_0Mr^n>R^n$, and then $\varepsilon_0>0$ with
\[
 \tfrac{16}{7}\varepsilon_0\leq\min\{(2M)^{-1},M^{-2}\}.
\]
The desired bound holds unless both
\[
 \|h_1-h_2\|_\infty\leq\kappa\delta s,
 \qquad\langle(b-a)_+\rangle_\omega\leq\varepsilon_0s.
\]
Assume both and put $t=\kappa s\leq1/4$. On $\Omega=\{(1-t)b-a>s/2\}$, we have $b-a>7s/16$, so its relative volume is at most $16\varepsilon_0/7$.

Take a smooth $\zeta:\mathbb R\to[0,1]$, zero on $(-\infty,1/2]$ and one on $[3/4,\infty)$, and set
\[
 \rho=b_0+M\zeta\bigl(((1-t)b-a)/s\bigr),\qquad
 b_0=1-M\bigl\langle\zeta\bigl(((1-t)b-a)/s\bigr)\bigr\rangle_\omega.
\]
The support lies in $\Omega$, so $1/2\leq b_0\leq1$ and $\langle\rho^2\rangle_\alpha\leq2+2M^2(16\varepsilon_0/7)\leq4$. Averages for $\alpha$ and $\omega$ agree. Thus the preceding auxiliary solution has the same fixed $c_0,C_\psi$.

At a point where $b-a=s$, the function $\Phi=(1-t)b-a+t\psi/2$ is at least $7s/8$. At its maximum $x_0$, $(1-t)b-a\geq7s/8$, so $\rho(x_0)\geq M$. The second derivative test gives, at $x_0$,
\[
 P=\alpha+\ddc\bigl((a-(1-t)b)/t\bigr)
 \geq\tfrac12(\alpha+\ddc\psi)+\tfrac12\alpha>0.
\]
The identity
\[
 \theta+\ddc a
 =t\bigl((\theta-\alpha)+P\bigr)+(1-t)(\theta+\ddc b)
\]
and concavity imply
\[
 F\bigl(\omega^{-1}((\theta-\alpha)+P)\bigr)
 \leq h_2+(h_1-h_2)/t\leq H+\delta.
\]
Both matrices in the convex combination are admissible. By \cref{lem:strict-gap-matrix-bound}, $P\leq R\omega$. But the auxiliary equation gives
\[
 R^n\geq P^n/\omega^n
 \geq2^{-n}c\rho r^n\geq2^{-n}c_0Mr^n>R^n,
\]
a contradiction. One of the two smallness inequalities must fail, giving the one-sided bound for $s$. Interchanging the solutions proves the theorem.
\end{proof}

\Needspace{9\baselineskip}
\begin{corollary}\label{cor:concave-uniform-compactness}\index{Compactness!concave equation solutions}
If the positive right-hand sides form an equicontinuous family with the same fixed data, the normalized smooth admissible solutions are precompact in $C^0(X)$ and have a common modulus of continuity.
\end{corollary}
\begin{proof}
Admissibility gives $\operatorname{tr}_\omega\ddc u\geq-C$. Fix the smooth positive invariant measure for this fixed elliptic operator, furnished by its Gauduchon adjoint. Integrating the product rule for $u^2$, using $0\leq u\leq C$ from \cref{prop:strict-gap-potential-bound}, bounds $\int_X|\mathrm du|_\omega^2\omega^n$ uniformly. The invariant density is uniformly comparable to $\omega^n$. Rellich compactness therefore gives $L^1$ precompactness.

From any sequence, Arzel\`a--Ascoli gives a uniformly convergent subsequence of right-hand sides. Extract a further $L^1$-Cauchy subsequence of solutions. \Cref{thm:concave-linear-stability} makes it uniformly Cauchy, proving $C^0$ precompactness. A finite uniform $\varepsilon/3$-net, together with the continuity of its finitely many functions, supplies a common modulus of continuity.
\end{proof}
